\documentclass{article}

\usepackage{booktabs}
\usepackage{tabularx}

\title{Development Plan\\\progname}

\author{\authname}

\date{}

\input{../Comments.text}
\input{../Common.text}

\begin{document}

\maketitle

\begin{table}[hp]
\caption{Revision History} \label{TblRevisionHistory}
\begin{tabularx}{\textwidth}{llX}
\toprule
\textbf{Date} & \textbf{Developer(s)} & \textbf{Change}\\
\midrule
Date1 & Name(s) & Description of changes\\
Date2 & Name(s) & Description of changes\\
... & ... & ...\\
\bottomrule
\end{tabularx}
\end{table}

\newpage{}

\wss{Put your introductory blurb here.  Often the blurb is a brief roadmap of
what is contained in the report.}

\wss{Additional information on the development plan can be found in the
\href{https://gitlab.cas.mcmaster.ca/courses/capstone/-/blob/main/Lectures/L02b_POCAndDevPlan/POCAndDevPlan.pdf?ref_type=heads}
{lecture slides}.}

\section{Confidential Information?}

\wss{State whether your project has confidential information from industry, or
not.  If there is confidential information, point to the agreement you have in
place.}

\wss{For most teams this section will just state that there is no confidential
information to protect.}
\section{IP to Protect}

\wss{State whether there is IP to protect.  If there is, point to the agreement.
All students who are working on a project that requires an IP agreement are also
required to sign the ``Intellectual Property Guide Acknowledgement.''}

\section{Copyright License}

\wss{What copyright license is your team adopting.  Point to the license in your
repo.}

\section{Team Meeting Plan}

\wss{This section is for meetings between the team members, and meetings between 
the team and their supervisor/advisor (as appropriate).}

\wss{How often will the team meet? where? when?}

\wss{If the meeting is a physical location (not virtual), out of an abundance of
caution for safety reasons you shouldn't put the location online}

\wss{To help schedule meetings, you can use the lecture/tutorial time when we do 
not have anything else scheduled.}

\wss{How often will you meet with your supervisor/advisor?  when?  where?}

\wss{Will meetings be virtual?  At least some meetings should likely be
in-person.}

\wss{How will the meetings be structured?  There should be a chair for all meetings.  
There should be an agenda for all meetings.}

\section{Team Communication Plan}

\wss{Issues on GitHub should be part of your communication plan.}

\section{Team Member Roles}

\wss{You should identify the types of roles you anticipate, like notetaker,
leader, meeting chair, reviewer.  Assigning specific people to those roles is
not necessary at this stage.  In a student team the role of the individuals will
likely change throughout the year.}

\section{Workflow Plan}

\begin{itemize}
	\item How will you be using git, including branches, pull request, etc.?
	\item How will you be managing issues, including template issues, issue
	classification, etc.?
  \item Use of CI/CD
\end{itemize}

\section{Project Decomposition and Scheduling}

\begin{itemize}
  \item How will you be using GitHub projects?
  \item Include a link to your GitHub project
\end{itemize}

\wss{How will the project be scheduled?  This is the big picture schedule, not
details. You will need to reproduce information that is in the course outline
for deadlines.}

\section{Proof of Concept Demonstration Plan}

Based on our project's scope and requirements, we have identified several
potential risks that we plan to address. These include:

\begin{itemize}
\item \textbf{Signal acquisition:} The system relies on us being able to
  accurately capture and interpret the Doppler ultrasound signal that the Vivid Q
  ultrasound machine outputs. We currently expect this to be a pair of
  quadrature-demodulated channels, which must be sampled simultaneously to
  identify flow directions and velocities. If we are unable to get a clean
  signal, or if there are unexpected voltage and noise levels, then we may need
  to add some additional processing before being able to process and interpret
  it as intended. This would have an impact on all further stages of the
  project.
\item \textbf{Mean Blood Velocity (MBV) estimation:} The MBV estimation process
  has several steps. These include combining the I and Q channels into a single
  complex signal, filtering out low-frequency, high-amplitude artifacts from
  vessel wall motion, and then processing using Fourier transforms to split the
  signal into forward and reverse flows. We will then need to estimate frequency
  shifts from the signal spectrum and calculate the velocities from there. With
  this, there are risks of noise, filter settings, aliasing, or spectral
  broadening, which could occur at any phase and would make our velocity output
  inaccurate. Because of this, we will be looking into several methods of doing
  these steps and will be prototyping before implementation to ensure that we
  have an accurate way of doing all of this.
\item \textbf{Real-time performance and hardware:} Because the system needs to
  work in real-time, we will need to make sure that we account for latency in
  our signal processing and the hardware we use. For a POC demo we expect to be
  running on a computer with recorded and simulated test data rather than
  low-level real-time hardware, as we will still be prototyping, but this is
  something we will need to account for in our design for the final system.
\item \textbf{Output validation:} Because the output of this system is going to
  be used for medical research, we need to ensure that it is accurate, both in
  measurement and in timing. The lab's calibration procedures rely on this
  signal, so we will need to test thoroughly against their old test and
  calibration data to make sure that our solution is both precise and accurate.
\end{itemize}

These risks may change as we learn more about the equipment and the signal. Our
proof of concept will demonstrate a minimal pipeline running on a computer,
which will go from the Doppler I/Q input to an output voltage signal. This will
allow us to validate our signal processing approach, while leaving the
implementation of low-level hardware development and real-time digital signal
processing for later in the development process. This proof of concept will
confirm that our signal processing steps are accurate and produce the desired
output, and we will ensure that this is correct before moving on to the
real-time and hardware aspects.

\section{Expected Technology}

Our technology choices listed below are based on what we currently know about
the project, and we may see them change as we learn more about the signal and
equipment in the following weeks. Our plan is to develop this system in two
stages. The first will be some high-level processing to ensure that our
approach is correct and to develop our proof of concept. The next stage will
involve moving the validated process on to hardware to run in real-time. Below
are the technologies we expect to use during both of these stages:

\begin{itemize}
\item \textbf{Programming Languages:} We plan to use Python for our initial
  prototyping, analysis, and proof of concept, as it allows us to easily test
  different signal processing methods and procedures, and visualize our
  results. The final product, which runs low-level on its own hardware, will be
  written in C or C++, since microcontrollers like the Teensy cannot run Python
  and we need efficient code to process the signal in real-time, consistent
  with the project requirements.
\item \textbf{Libraries:} During the POC phase, we expect to use a number of
  Python libraries such as NumPy and SciPy for signal processing, filtering, and
  transforms. Matplotlib will be used to visualize our traces and signals
  before, during, and after each step of the signal processing. For the
  real-time system, we will use whichever libraries are compatible with the
  hardware we select. One current proposed idea is to use a Teensy 4.1 board
  with an external DAC. In this case, we would use the Teensy Audio Library to
  capture the stereo input, along with ARM's CMSIS-DSP library for the complex
  transforms and filtering.
\item \textbf{Hardware:} We expect to use some analog circuitry to connect to
  the Doppler input and filter out interference before it is digitized, as well
  as a commercial off-the-shelf board that supports DSP, such as a Teensy 4.1.
  With this, we may use an external digital-to-analog converter (DAC), as well
  as adapters to output to the physical wired connections that the PowerLab
  8/30 device takes as input for LabChart. Further details of the hardware we
  will use are still to be determined, as we are currently exploring our
  options.
\item \textbf{Linting and Formatting:} We plan to use Pylint to lint our code
  during the proof of concept phase, and Black for Python formatting. Later on,
  in the final development stage when we shift to C or C++, we expect to use
  Clang-Format, Clang-Tidy, and Cppcheck for our linting and formatting.
\item \textbf{Testing:} We will use Pytest for our Python development and a
  framework like Unity for C/C++ testing. Because our project also involves the
  use of hardware, we will use LabChart output and an oscilloscope to verify the
  final system output.
\item \textbf{Continuous Integration:} We will use GitHub Actions to run the
  linter and unit tests on every pull request. Since code that depends on
  physical hardware cannot run in CI, we will keep the processing logic
  separate from the hardware input and output code.
\item \textbf{Version Control:} We will be using Git and GitHub for version
  control, along with GitHub Issues and GitHub Projects for task tracking.
\item \textbf{Development Environment:} We will use Visual Studio Code as our
  main IDE for development as well as for our \LaTeX{} documentation.
\item \textbf{Laboratory Equipment:} We will be using the Vascular Dynamics
  Lab's hardware for some testing and integration, which includes a GE Vivid Q
  ultrasound, a Multigon 500M Transcranial Doppler, a PowerLab 8/30 unit, and
  the LabChart software for signal input, recording, comparison, and
  validation.
\item \textbf{Components Implemented Ourselves:} While we will be using
  libraries for standard building blocks like Fourier transforms and filters in
  both Python and C/C++, we plan to implement project-specific parts of the
  signal processing ourselves. This may include arranging the input transform
  and separating forward and reverse flow, designing the initial filter for the
  lab's measurements, estimating the mean frequency from the spectrum,
  converting the Doppler shift to velocity, and calibrating the final output.
\end{itemize}

\section{Use of AI and LLMs}

\wss{This section is a continuation of the previous section.  The role of LLMs
in your project as important enough to warrant their own section.  
How will your team be using LLMs for code and documentation development?  
What tools will you be using?  How will you verify the output of your LLMs?}

\section{Coding Standard}

\wss{What coding standard will you adopt?}

\newpage{}

\section{Appendix --- Generative AI Use Disclosure}

\wss{ChatGPT (version 5) was used to brainstorm ideas for the template for this disclose section.}

\wss{This section is about the use of AI to write this document, not your
planned use of AI for your project.  You discuss that topic in one of the
sections in the body of this report.}

\subsection{AI Tools Used}

\wss{Give the name of the AI tool(s) used and the version.}

\subsection{How and Where Used}

\wss{Explain what the tool(s) were used for.  Examples include: brainstorming,
explaining a technical concept, reviewing the document, generating or modifying
text, generating or modifying diagrams, identifying errors or omissions or
inconsistencies.}

\wss{For each of the uses, identify which parts of the document were affected.}

\wss{You do not need to report every interaction, but you should disclose the 
uses that materially contributed to your work}

\subsection{Verification of AI-Generated Content}

\wss{\begin{itemize}
    \item Describe how the team checked AI-generated or AI-assisted material
    before including it in the document.
    \item In particular, verify factual claims, technical assertions,
    requirements, calculations, references, and other information for which an
    AI system may produce incorrect or fabricated results.
    \item \textbf{The team is responsible for the accuracy, completeness,
consistency, and appropriateness of the submitted document, regardless of
whether material was generated by a human or an AI system.}
\end{itemize}}

\newpage{}

\section*{Appendix --- Reflection}

\wss{Not required for CAS 741}

\input{../Reflection.text}

\begin{enumerate}
    \item Why is it important to create a development plan prior to starting the
    project?
    \item In your opinion, what are the advantages and disadvantages of using
    CI/CD?
    \item What disagreements did your group have in this deliverable, if any,
    and how did you resolve them?
\end{enumerate}

\newpage{}

\section*{Appendix --- Team Charter}

\wss{borrows from
\href{https://engineering.up.edu/industry_partnerships/files/team-charter.pdf}
{University of Portland Team Charter}}

\subsection*{External Goals}

\wss{What are your team's external goals for this project? These are not the
goals related to the functionality or quality fo the project.  These are the
goals on what the team wishes to achieve with the project.  Potential goals are
to win a prize at the Capstone EXPO, or to have something to talk about in
interviews, or to get an A+, etc.}

\subsection*{Attendance}

\subsubsection*{Expectations}

\wss{What are your team's expectations regarding meeting attendance (being on
time, leaving early, missing meetings, etc.)?}

\subsubsection*{Acceptable Excuse}

\wss{What constitutes an acceptable excuse for missing a meeting or a deadline?
What types of excuses will not be considered acceptable?}

\subsubsection*{In Case of Emergency}

\wss{What process will team members follow if they have an emergency and cannot
attend a team meeting or complete their individual work promised for a team
deliverable?}

\subsection*{Accountability and Teamwork}

\subsubsection*{Quality} 

\wss{What are your team's expectations regarding the quality
of team members' preparation for team meetings and the quality of the
deliverables that members bring to the team?}

\subsubsection*{Attitude}

\wss{What are your team's expectations regarding team members' ideas,
interactions with the team, cooperation, attitudes, and anything else regarding
team member contributions?  Do you want to introduce a code of conduct?  Do you
want a conflict resolution plan?  Can adopt existing codes of conduct.}

\subsubsection*{Stay on Track}

\wss{What methods will be used to keep the team on track? How will your team
ensure that members contribute as expected to the team and that the team
performs as expected? How will your team reward members who do well and manage
members whose performance is below expectations?  What are the consequences for
someone not contributing their fair share?}

\wss{You may wish to use the project management metrics collected for the TA and
instructor for this.}

\wss{You can set target metrics for attendance, commits, etc.  What are the
consequences if someone doesn't hit their targets?  Do they need to bring the
coffee to the next team meeting?  Does the team need to make an appointment with
their TA, or the instructor?  Are there incentives for reaching targets early?}

\subsubsection*{Team Building}

\wss{How will you build team cohesion (fun time, group rituals, etc.)? }

\subsubsection*{Decision Making} 

\wss{How will you make decisions in your group? Consensus?  Vote? How will you
handle disagreements? }

\end{document}